% Vietnamese translation for OI Public Library
% Source-File: IOI中国国家候选队论文/2003/刘一鸣/刘一鸣.ppt
\documentclass[11pt]{article}
\usepackage{oipl-vn}

\title{Bản trình chiếu: Sắp thứ tự dữ liệu trong bài toán tìm kiếm}
\author{Liu Yiming}
\date{2003}

\begin{document}
\maketitle

\origtitle{Optimization ideas for a class of search problems: ordering data}
\sourcepath{IOI China candidate team papers / 2003 / Liu Yiming / Liu Yiming.ppt}

\section{Tư tưởng mở đầu}

Bài trình bày nói về một kiểu tối ưu hóa rất thực dụng cho tìm kiếm: biến dữ
liệu lộn xộn thành dữ liệu có trật tự trước khi hoặc trong khi tìm kiếm. Khi
dữ liệu có cấu trúc, ta thường gặp nghiệm tốt sớm hơn, cắt tỉa mạnh hơn, và
tránh xét nhiều trạng thái đẳng cấu.

Ví dụ nhập môn là bài xếp hàng vào container. Nếu mục tiêu là lấp đầy đúng
thể tích, thứ tự xét hàng hóa ảnh hưởng trực tiếp đến nghiệm đầu tiên mà DFS
tìm được. Sắp xếp hàng hóa trước khi tìm giúp các cận tối ưu phát huy sớm
hơn, dù bản thân bước sắp xếp chỉ là một thay đổi nhỏ trong chương trình.

\section{Hai kiểu sắp thứ tự dữ liệu}

Slide chia phương pháp thành hai nhóm:
\begin{itemize}
  \item sắp thứ tự trong tiền xử lý, tức là chuẩn hóa dữ liệu ngay sau khi đọc
  vào rồi mới chạy thuật toán chính;
  \item sắp thứ tự trong thời gian thực, tức là mỗi trạng thái sinh ra trong
  quá trình tìm kiếm được chuẩn hóa trước khi kiểm tra hoặc lưu trữ.
\end{itemize}

\section{Tiền xử lý: Blocks}

Ví dụ chính cho tiền xử lý là bài \textit{Blocks} của IOI 2000. Bài toán hình
học ba chiều ban đầu được đổi thành bài toán tập hợp một chiều. Ta sắp xếp
các khối lập phương nhỏ trong cấu hình, đánh số chúng từ \(1\) đến \(v\), rồi
dùng tập chỉ số để biểu diễn cấu hình.

Mỗi khối có thể đặt vào cũng được biểu diễn bằng tập các chỉ số ô mà nó phủ.
Khi lấy đi một khối
\[
  \{3,6,7,9\}
\]
khỏi cấu hình \([1,10]\), phần còn lại là
\[
  \{1,2,4,5,8,10\}.
\]
Muốn kiểm tra một khối khác, chẳng hạn
\[
  \{4,5,7,8\},
\]
ta chỉ cần kiểm tra quan hệ tập con; vì \(7\) đã bị lấy đi, khối này không
thể đặt tiếp. Tương tự, hai khối xung đột khi giao của hai tập biểu diễn khác
rỗng:
\[
  \{3,6,7,9\}\cap\{4,5,7,8\}=\{7\}.
\]

Sau bước chuẩn hóa này, phần tìm kiếm chính không còn phải xử lý nhiều chi
tiết hình học ba chiều. Đây là lợi ích lớn nhất của sắp thứ tự trong tiền xử
lý: làm mô hình bài toán gọn hơn trước khi DFS bắt đầu.

\section{Thời gian thực: biểu diễn nhỏ nhất}

Nhóm kỹ thuật thứ hai xử lý trạng thái ngay khi nó được sinh ra. Công cụ tiêu
biểu là \term{biểu diễn nhỏ nhất}: trong một lớp trạng thái đẳng cấu, chỉ giữ
một đại diện chuẩn, thường là đại diện nhỏ nhất theo thứ tự từ điển.

Nếu dùng cách truyền thống, khi sinh được một trạng thái hợp lệ ta phải sinh
mọi trạng thái đối xứng của nó để xem đã lưu trạng thái nào tương đương chưa.
Với biểu diễn nhỏ nhất, ta chuyển trạng thái hiện tại về đại diện chuẩn rồi
chỉ kiểm tra đại diện đó. Tính chất quan trọng là tính duy nhất: mỗi lớp đẳng
cấu có đúng một biểu diễn nhỏ nhất.

\section{Ví dụ N quân hậu}

Với bài N quân hậu, một bàn cờ được biểu diễn bằng dãy
\[
  a_1,a_2,\ldots,a_n,
\]
trong đó \(a_i\) là cột của quân hậu ở hàng \(i\). Các phép đối xứng của bàn
cờ như lật ngang, lật dọc, lật qua đường chéo và quay \(90^\circ\),
\(180^\circ\), \(270^\circ\) đều biến dãy này thành một dãy khác. Với
\(n=5\), các công thức kiểu đảo thứ tự hoặc thay \(a_i\) bằng \(6-a_i\) chính
là cách biểu diễn những phép đối xứng đó trên dãy.

Trong quá trình quay lui, nếu tiền tố hiện tại chắc chắn không thể phát triển
thành biểu diễn nhỏ nhất của lớp đối xứng, ta có thể cắt nhánh ngay. Nhờ vậy,
không cần lưu toàn bộ nghiệm đã gặp để so sánh đẳng cấu; ta chỉ liệt kê các
đại diện chuẩn. Cách này vừa giảm thời gian kiểm tra, vừa giảm bộ nhớ.

\section{So sánh và kết luận}

Sắp thứ tự trong tiền xử lý thường đơn giản và nhanh, nhưng có thể dùng nhiều
bộ nhớ vì phải chuẩn bị cấu trúc dữ liệu trước. Sắp thứ tự trong thời gian
thực linh hoạt hơn và tiết kiệm bộ nhớ hơn, nhưng có thể tốn chi phí lặp lại
ở nhiều nút của cây tìm kiếm. Hai phương pháp bổ sung cho nhau.

Thông điệp cuối của slide là không nên xem sắp thứ tự dữ liệu như một mẹo
nhỏ. Với nhiều bài tìm kiếm, chỉ vài dòng chuẩn hóa dữ liệu có thể biến không
gian tìm kiếm rối thành một mô hình gọn, giúp các cắt tỉa khác hoạt động tốt
hơn và tránh xét lại các trạng thái đối xứng.

\end{document}
